\documentclass[UTF8]{ctexart}
\usepackage{amsmath,amssymb}
\usepackage{enumitem}
\setlist[itemize]{leftmargin=*, itemsep=0pt, topsep=2pt}
\usepackage[hidelinks]{hyperref}

\begin{document}
\section*{工作大纲}


\subsection*{0. 名称与定位（WS-OVVIS基准与协议化评测）}
WS-OVVIS（Weakly-supervised Open-Vocabulary Video Instance Segmentation）：在仅给clip-level正证据标签集合（positive label set）且集合系统性不完备的弱监督条件下学习开放词汇视频实例分割；并提供统一的缺失标签生成协议、unknown评测与鲁棒性曲线，形成可复现的对比基准。


\par\noindent\hrulefill\par


\subsection*{1. 任务设定}
输入监督：
\begin{itemize}
\item 每个视频段（video clip）的语义类别列表（clip-level label set，多标签集合）
\item clip-level label set $Y'$正确但不完备：$Y'\subseteq Y^*$（$Y^*$为该clip的真实类别集合），并按缺失协议构造（见4.1）；允许open set（存在$Y^*$之外的实例类别/测试时出现novel classes）
\begin{itemize}
\item Uniform Missing：对$Y^*$中每个真实类别以概率$p$被drop
\item Long-tail Missing：对长尾类设置更高drop概率（如随类频率反比）
\item Scenario/Domain Missing（可选）：按场景/域对部分类系统性缺失
\end{itemize}
\item 不使用任何像素级标注（mask/box/point），不使用tube-level paired training data（不使用实例级类别标注与实例级配对监督）
\item 不依赖任何密集分割大模型教师（如 SAM/Mask2Former）
\item 使用冻结的CLIP Text Encoder（默认，只用文本侧，不用CLIP image encoder）或LLaMA3（可选）生成类别文本原型
\end{itemize}
目标：
\begin{itemize}
\item 在上述弱监督条件下学习一个开放词汇的视频实例分割模型（OV-VIS）
\item 同时提供WS-OVVIS训练标签缺失协议、评测指标与baseline套件，使该设定可复现、可直接对比
\item 同时输出video instance的mask序列、track ID与类别预测，支持未见类别（novel classes）
\end{itemize}

\par\noindent\hrulefill\par


\subsection*{2. 核心技术路线（SeqFormer + set-to-track归因 + 闭环自训练）}
整体思路：
\begin{itemize}
\item 用无监督/弱监督tube提案提供class-agnostic伪监督，训练一个video-native的实例分割器获得稳定的track-level instance queries
\item 用clip-level label set进行多实例归因（set-to-track assignment），将集合监督转化为实例级语义监督
\item 用文本原型提供开放词汇类别空间，并在训练中将语义注入到track-level表征
\end{itemize}

\subsubsection*{2.1 伪tube生成与预处理}
\begin{itemize}
\item 使用VideoCutLER在每个video clip中生成class-agnostic pseudo instance mask tube（含tracklet）
\item 对tube做质量筛选与去重（时序一致性、面积稳定性、mask重叠/NMS、短轨过滤）
\item 生成训练所需的伪监督：每帧mask、tube/track ID、tube置信度权重
\item 说明：伪tube生成器是可替换部件（可在实验中替换为其它proposal/tracklet来源），WS-OVVIS协议不依赖特定实现。
\end{itemize}

\subsubsection*{2.2 video-native instance segmentor}
\begin{itemize}
\item 采用SeqFormer作为video instance segmentor
\item 训练目标（class-agnostic）：
\begin{itemize}
\item mask loss：用伪mask tube监督每个query的mask序列
\item association/ID：通过Hungarian匹配（基于 mask IoU/box IoU/时序一致性）将query与伪tube对齐
\item 训练时不使用实例类别GT，仅学习实例分割与跨帧一致的instance queries（video-level instance embedding）
\end{itemize}
\item 主线进展锚点（mainline closeout）：
\begin{itemize}
\item 基线全规模链路已完成到`model_final.pth`并完成同谱系eval恢复（见P60--P62 closeout链）。
\item method-active全规模链路已在同一rerun root完成终态checkpoint与post-train eval产物（见P65--P68 closeout链）。
\end{itemize}
\end{itemize}

\subsubsection*{2.3 track-level表征构建（归因输入）}
\begin{itemize}
\item 从SeqFormer获得每个video-level instance query的表征（作为track embedding）
\item 形成当前视频的candidate track feature set：$F = \{f_i\}，i=1 \dots N$
\item 为每个track维护置信度（来自mask质量、时序一致性、匹配稳定性）
\end{itemize}

\subsubsection*{2.4 文本原型构建}
\begin{itemize}
\item 冻结CLIP Text Encoder，输入类别名与提示模板（prompt ensemble，可选类别描述短句），得到文本原型集合$G=\{g_j\}$
\item 或使用LLaMA3（或其它纯文本encoder）编码类别名与描述，得到$G=\{g_j\}$
\end{itemize}

\subsubsection*{2.5 多实例归因（set-to-track assignment，核心模块）}
输入：
\begin{itemize}
\item track feature set $F=\{f_i\}$
\item 当前视频的label set $Y'$（正证据集合，只包含部分出现的类别；可能缺失真实出现类）
\item 文本原型$G=\{g_j\}$（全词表），并构建候选类别集合$\hat{Y}=Y'\cup TopK(F, G)$；归因列包含$\hat{Y}$ + bg + unk-fg（集合外前景）。
\end{itemize}
相似度矩阵：
\begin{itemize}
\item $S_{ij} = sim(f_i, g_j)$（点积或余弦，相似度温度$\tau$）
\end{itemize}
归因策略（主方法 + 备选）：
\begin{itemize}
\item 主方法：OT/Sinkhorn软分配得到$P$（track × label），在$\hat{Y}\cup\{bg,unk\}$上软分配得到$P$
\begin{itemize}
\item 对$y\in Y'$：加入coverage约束/损失（保证每个$Y'$中的类别被至少一些track解释）。
\item 对$c\in \hat{Y}\setminus Y'$：不施加排他性约束（允许 open set）。
\item 仍保留背景列，但新增unk-fg列承接“高objectness但不匹配”的track，避免压bg。
\end{itemize}
\item 备选：MIL（max/log-sum-exp聚合）或EM（软后验迭代）作为对照与消融
\end{itemize}
训练注入：
\begin{itemize}
\item 语义对齐损失：用$P$作为软标签监督track embedding的语义分布
\item 背景抑制：背景列概率用于降低噪声track的梯度贡献（loss weighting）
\item 时序语义一致性：同一track在不同帧/不同片段增强下的语义分布一致（KL/InfoNCE/一致性正则）
\item Coverage loss（正证据覆盖）：对每个$y\in Y'$，使用Noisy-OR/top-r pooling的bag-level约束，保证$Y'$中类别在该clip内可被解释（“必须覆盖正证据集合$Y'$”，但不限制只能是$Y'$）。
\item FG-not-BG正则（防背景塌缩）：对高置信track（来自mask质量/时序一致性）抑制其被分配到bg，促使其进入unk-fg或open类。
\end{itemize}
当前主线证据已确认该模块可在method-active全规模训练中以non-zero模式被配置并运行到终态（P65--P68）；但最终跨方法标量收益表仍待整理（见4.2/4.4中的TODO）。

\subsubsection*{2.6 闭环自训练与伪tube纠错}
\begin{itemize}
\item 以当前模型预测对伪tube做refinement：
\begin{itemize}
\item 合并碎片tracklet、剔除背景tube、修正mask边界与漂移
\end{itemize}
\item 进行1–2轮迭代训练，逐步降低对初始伪tube的依赖
\item 训练课程：先用高置信 tube，再逐步引入中低置信tube
\item Label-set expansion：将unk-fg中高置信track的Top-1/Top-2 open类加入下一轮该clip的$\hat{Y}$（仅作为候选列参与归因，不回写为“真标签”），逐轮提升集合外类别可分配性。
\end{itemize}
当前已闭合的Stage D有界机制链（D23--D25）支持“guard/fallback-guard可稳定阻断round2回退”的机制结论，但该证据为有界cohort结论，不替代完整大规模消融。

\subsubsection*{2.7 推理（开放词汇VIS）}
\begin{itemize}
\item 对每个视频clip，SeqFormer输出instance masks + video-level instance queries
\item 计算query embedding与全词表文本原型的相似度，获得类别分布
\item 输出：mask序列 + track ID + open-vocab类别（支持novel类）
\end{itemize}

\par\noindent\hrulefill\par


\subsection*{3. 创新点与贡献}
\begin{itemize}
\item 提出WS-OVVIS：弱监督（clip-level positive label set）+ 标签不完备 + open-vocab VIS的新任务与可复现实验协议（缺失标签生成、评测脚本、baseline套件与splits）
\item 系统化clip-level label set监督下的弱监督开放词汇视频实例分割设定（不使用tube-level paired training data，不使用像素级监督）
\item 提出set-to-track多实例归因框架，将集合级语义监督转化为track-level语义监督，显式建模背景与噪声track，并与video-native instance queries联合优化
\item 构建video-native学习闭环：SeqFormer track-level表征 + 时序语义一致性 + tracklet处理 + 自训练refinement，提高在噪声tube下的稳定性与上限
\item 第一个把OV-VIS推到clip-level label-set supervision的工作
\item 将clip-level label set从“封闭世界类别空间”重述为正证据覆盖约束，提出open-world set-to-track attribution（含unk-fg显式建模与 coverage loss），在保证覆盖Y的同时允许实例预测超出Y，避免集合外前景被压为背景。
\end{itemize}

\par\noindent\hrulefill\par


\subsection*{4. 核心实验设置}

\subsubsection*{4.1 数据集与评测协议}
\begin{itemize}
\item 主数据集：LV-VIS（4,828 videos / 1,196 categories；train/val/test=3,083/837/980；base/novel=641/555），按base/novel划分报告AP、$AP_{base}$、$AP_{novel}$
\item WS-OVVIS训练监督构造：对训练集每个clip由其真实类别集合$Y^*$构造不完备集合$Y'$（仅训练时使用$Y'$；val/test评测仍使用原始GT，不做缺失）
\begin{itemize}
\item Missing rate：$p \in \{0.0, 0.2, 0.4, 0.6\}$（至少3个点，用于鲁棒性曲线）
\item 缺失模式：Uniform Missing / Long-tail Missing（可选再加Scenario/Domain Missing）
\end{itemize}
\item 训练/调参/汇报口径：
\begin{itemize}
\item 训练仅使用 LV-VIS train（3,083）；伪tube可对 train/val/test 全集推理生成，但不使用其GT。
\item 超参选择与消融在 LV-VIS val（837）完成；最终主结果在 LV-VIS test（980）汇报（如test需服务器提交，则同时在val给出可复现实验表）。
\end{itemize}
\item 泛化评测：在LV-VIS训练后直接迁移到YouTube-VIS19/21、OVIS、BURST的val集做评测（不在目标数据集上finetune，不使用其像素标注训练）
\item 训练可用信息：视频帧 + clip-level label set + VideoCutLER伪tube/mask及自训练更新后的伪tube
\item 指标：VIS标准AP/AP50/AP75（AP为IoU阈值0.50:0.05:0.95平均；IoU按时空mask计算）；open-vocab分解$AP_{base}/AP_{novel}$
\item WS-OVVIS专用指标（用于“标签不完备 + unknown”可衡量）：
\begin{itemize}
\item Missing-Label Robustness Curve：mAP vs missing rate $p$，并报告AURC（Area Under Robustness Curve）
\item Set-Coverage Recall（SCR）：对每个clip衡量$Y'$中每个类别是否被至少一个预测实例“解释”（Noisy-OR/top-r pooling定义）
\item Unknown-handling诊断（可选）：对$Y^*\setminus Y'$对应实例统计class-agnostic跟踪/分割质量，验证unk-fg是否承接集合外前景而非压bg
\end{itemize}
\end{itemize}


\subsubsection*{4.1.1 训练流程与公平性口径（Stage A--D）}
\begin{itemize}
\item Stage A（伪tube生成）：用VideoCutLER在LV-VIS train/val/test推理生成class-agnostic pseudo mask tubes；固定checkpoint与过滤/去重阈值；所有方法与消融共享同一份伪tube。
\item Stage B（video-native class-agnostic 分割器）：用伪tube监督训练SeqFormer（固定$T$帧采样与Hungarian匹配；BCE+Dice等mask损失；tube置信度re-weight），得到稳定的track features $F=\{f_i\}$与mask序列。
\item Stage C（open-world set-to-track 归因训练）：构建$\hat{Y}=Y'\cup TopK(F,G)$，计算相似度$S_{ij}$（温度$\tau$）；用OT/Sinkhorn得到$P$并加入coverage、bg、unk-fg；对照实现MIL/EM（共享相同$F,G,\hat{Y},\tau,TopK$与训练轮数）。
\item Stage D（闭环自训练）：用当前模型refine/更新伪tube并再训练（0/1/2轮）；课程学习从高置信到低置信；报告不同轮数对AP/APn与鲁棒性的影响。
\item 主线完成态（截至本draft）：baseline全规模链路已闭合（train终态+eval恢复产物）；method-active全规模链路已闭合（同根continuation终态+post-train eval产物）；Stage D机制链完成有界证据封存（3-case cohort）。
\end{itemize}

\subsubsection*{4.1.2 Stage D机制结果（bounded round2 continuation）}
Stage D在三个固定Stage C summary上执行bounded round2 continuation，对比unguarded continued additions与guarded/fallback-guarded执行。该协议仅检验一个机制问题：round2继续加候选是否会触发loss回退，以及guard逻辑是否能稳定该行为。

在三个summary中，unguarded分支均接受同一round2 proposal（`909002`），候选数从7增至8，并在round2达到loss `5.571925640106201`。对应地，guarded/fallback-guarded分支在同一proposal上给出reject，候选数保持7，round2 loss保持在`4.822475433349609`。这在当前cohort形成`3/3`回退阻断，mean round2 loss delta（guarded - unguarded）为`-0.7494502067565918`。

就ws指标而言，round2 AURC在两种条件下均为`0.5833333333333333`，mean delta为`0.0`，因此当前证据支持damage control而非ws AURC uplift。Route A keep-one变体（`accept_top1_only_under_guard_v1`）在同一cohort下为负结果：`3/3`案例相对guarded baseline回退，且无AURC提升。

上述结论为有界机制证据，可支持主文机制叙述；但其外推范围受限于当前cohort规模与协议边界。

\subsubsection*{4.2 必须对比的baseline}
\begin{itemize}
\item 公平性口径：所有baseline与消融共享同一份VideoCutLER伪tube（同checkpoint与过滤阈值），并在相同backbone/训练轮数/词表上对比。
\item VideoCutLER直接输出 + 语义贴标（CLIP/LLaMA3文本原型），不训练分割器
\item SeqFormer class-agnostic（仅伪tube监督）+ 推理阶段语义贴标（不做归因训练）
\item 帧级方法（MinVIS 风格）：按图像方式训练/推理，通过跨帧匹配构建 track，再语义贴标（用于反驳“视频不必要”）
\item 无归因版本：仅做表示对齐或MIL聚合，不做set-to-track分配（用于反驳“只是视频版对齐”）
\item 协议对照：将训练监督从$Y'$替换为“错误标签集合/更强缺失/随机标签集合”，验证方法对噪声与缺失的鲁棒性，并排除“刚好吃到干净标签”的解释
\item TODO(mainline-pack): 当前仅完成baseline与method-active闭环产物封存；统一口径的baseline对比标量表（AP/AP50/AP75及base/novel拆分）待后续提取与汇总。
\end{itemize}

\subsubsection*{4.3 上界对照（强监督OV-VIS）}
\begin{itemize}
\item OV2Seg、CLIP-VIS、OpenVIS/InstFormer（作为 fully-supervised/open-vocab 上界定位差距）
\item TODO(mainline-pack): 上界对照未在本次主线封存任务中新增可引用实验结果，仅保留对照框架与目标位置。
\end{itemize}

\subsubsection*{4.4 关键消融（必须项）}
\begin{itemize}
\item 归因模块：OT vs EM vs MIL；是否引入背景列；OT 约束（列覆盖/行容量/熵正则）
\item 归因粒度：track-level vs tube-level vs frame-level
\item video-native机制：时序语义一致性开/关；tracklet stitching开/关
\item 闭环自训练：0 轮 vs 1 轮 vs 2 轮 refinement；tube 过滤/课程学习策略
\item 缺失协议消融：Uniform vs Long-tail；不同missing rate $p$ 下的鲁棒性曲线与AURC变化
\item Two-track语义：
\begin{itemize}
\begin{itemize}
\item CLIP Text（默认） vs 无CLIP（LLaMA3）
\item 类名only vs prompt ensemble
\end{itemize}
\end{itemize}
\item TODO(mainline-pack): 关键消融的完整结果矩阵（含missing-rate曲线与模块切换对照）尚未在当前封存包中形成可直接入文表格。
\end{itemize}

\subsubsection*{4.5 定性与诊断分析}
\begin{itemize}
\item 归因矩阵$P$的可视化（track × label），展示训练早期到收敛的变化
\item 自训练前后伪tube与预测mask的改善案例
\item 失败模式分析：遮挡交换、同类多实例混淆、小目标与长尾类别
\item missing-label失败模式：当$Y'$丢失关键类时，模型是否错误压bg、是否由unk-fg承接、coverage是否误触发导致假阳性
\item TODO(mainline-pack): 该部分所需的主文/附录可视化面板仍待独立打包。
\end{itemize}

\par\noindent\hrulefill\par


\subsection*{5. 预期结果与结论组织方式}
\begin{itemize}
\item 最低达标性能（顶会主会底线，用于证明弱监督下仍具备open-vocab能力）：LV-VIS上达到test AP$\ge 11.0$且$AP_{novel}\ge 9.0$；同时val AP$\ge 14.0$且$AP_{novel}\ge 12.0$。参考强监督OV2Seg（R50）：val AP=14.2/$AP_{novel}$=11.9、test AP=11.4/$AP_{novel}$=8.9；参考Detic-XMem（R50）：val AP=8.8/$AP_{novel}$=5.4、test AP=7.7/$AP_{novel}$=3.6；上界参考CLIP-VIS（R50，val）：AP=18.8/$AP_{novel}$=22.8。
\item 泛化最低线：在不finetune条件下，YouTube-VIS2019 novel mAPn$\ge 10$、OVIS mAP$\ge 11$（参考OV2Seg：YTVIS2019 mAPn=11.1，OVIS mAP=11.2；BURST mAP=3.7）。
\item 主结果：默认使用 CLIP Text Encoder（冻结）作为语义原型，展示最强 open-vocab 性能
\item 框架独立性：无CLIP版本作为消融/附录结果，证明set-to-track归因与video-native闭环是核心贡献
\item 公平性声明：训练阶段不使用任何tube-level paired training data与像素标注，实例级监督仅来自伪tube与自训练闭环
\item 复现与发布：提供WS-OVVIS缺失标签生成脚本、评测代码（AURC/SCR）、baseline实现与splits，便于社区直接对比
\end{itemize}

\end{document}
